\section*{Chapter \ref{ch:b3:ricardian-trade} --- Ricardian Trade in Full}

\begin{solution}{exo:b3:ricardian-trade:1}
Negative in every year, as the law predicts. % figdata: japan
It is closest to zero in 1868, at $-0.18$ million ryō. The prediction is about
the sign, not the size; and 1868 was a year in which Japan ran a surplus on
current account, so the negative sign there owes nothing to borrowing.
% ledger: b3rt:bb04.surplus
\end{solution}

\begin{solution}{exo:b3:ricardian-trade:2}
$A = 10/2 = 5$, $9/3 = 3$, $6/4 = 1.5$, $6/6 = 1$, $3/5 = 0.6$: the chain is goods 1, 2,
3, 4, 5 in that order. At $\omega = 1.2$ Home makes goods 1, 2 and 3; at $\omega = 3.5$,
good 1 only.
\end{solution}

\begin{solution}{exo:b3:ricardian-trade:3}
\begin{enumerate}
\item Brenland $4/2 = 2$ units of grain per unit of cloth; Kestria $6/12 = 0.5$.
\item $p_{\text{cloth}}/p_{\text{grain}} = w^* \times 6/(w \times 2) = 6/(2 \times 2) =
1.5$, between the two \omterm{def:b3:ricardian-trade:autarky}{autarky prices}. Brenland exports grain: its \omterm{def:b3:ricardian-trade:terms-of-trade}{terms of trade}
are $1/1.5 = 0.67$ units of cloth per unit of grain, against 0.5 in autarky.
\end{enumerate}
\end{solution}

\begin{solution}{exo:b3:ricardian-trade:4}
Japan exporting goods whose prices had been high before 1859 relative to world
prices and importing goods that had been cheap: net trade valued at \omterm{def:b3:ricardian-trade:autarky}{autarky
prices} would then be positive.
\end{solution}

\begin{solution}{exo:b3:ricardian-trade:5}
$B(z) = z/(1 - z)$. Setting $3 - 2z = z/(1 - z)$ gives $(3 - 2z)(1 - z) = z$, or $2z^2 - 6z
+ 3 = 0$, so $\bar z = (6 - \sqrt{12})/4 = 0.634$ (the other root exceeds 1) and $\omega = 3
- 2 \times 0.634 = 1.732$.
\end{solution}

\begin{solution}{exo:b3:ricardian-trade:6}
\begin{enumerate}
\item $A_{\text{grain}} = 6$, $A_{\text{cloth}} = 1.5$. With Brenland making grain only,
half of world spending goes to it: $wL = \tfrac12(wL + w^* \times 2L)$, so $\omega = 2$,
which lies between 1.5 and 6 as \cref{prop:b3:ricardian-trade:cut} requires. Cloth
now costs Brenland $w^* \times 6 = 3w$, against $4w$ in autarky; grain is unchanged.
With equal shares, real income rises by $(4/3)^{1/2} - 1 = 15.5~\%$.
\item $A_{\text{grain}} = 3$, still above 2: the equilibrium, cloth's price and
Brenland's real income are unchanged.
\item $A_{\text{grain}} = 1.5 = A_{\text{cloth}}$: the relative \omterm{def:b3:ricardian-trade:autarky}{autarky prices} are equal,
\omterm{def:b3:ricardian-trade:comparative-advantage}{comparative advantage} is gone and so is the reason to trade. Brenland's real
income falls back to autarky, by $1 - 1/1.155 = 13.4~\%$. The gain abroad was in
the good Brenland exports, not uniform, and
\cref{prop:b3:ricardian-trade:abroad} does not apply.
\end{enumerate}
\end{solution}

\begin{solution}{exo:b3:ricardian-trade:7}
$\kappa\omega = 1.56$ and $\omega/\kappa = 1.083$. Home exports goods 1 and 2 ($A = 5$ and
3), above 1.56; imports goods 4 and 5 ($A = 1$ and 0.6), below 1.083; and makes good 3
($A = 1.5$), the only one inside the band, for itself.
\end{solution}

\begin{solution}{exo:b3:ricardian-trade:8}
The tested solution is \code{ex08_crops.py}, in the chapter's \code{exercises/}
folder. Wheat against maize: 73 countries, correlation $-0.30$; wheat against
barley: 77 countries, 0.22. % output: ex08_crops
Wheat and maize grow in different climates, and in a dry country maize is
often confined to a few irrigated fields with high yields, so actual yields
are strongly selected; wheat and barley grow on similar land in similar seasons,
and the selection may matter less. Actual yields answer neither question
cleanly, which is why the published test used potential yields.
\end{solution}

\begin{solution}{exo:b3:ricardian-trade:9}
$0.8^{-1/5} - 1 = 4.6~\%$; $0.8^{-1/10} - 1 = 2.3~\%$. $\lambda$ is measured from trade and
output data with modest error; the elasticity is estimated, and halving it
doubles the answer, so it is the input to know better.
\end{solution}

\begin{solution}{exo:b3:ricardian-trade:10}
Home's workers make the goods it exports, $[0, z_E]$, bought by both countries,
and the \omterm{def:b3:ricardian-trade:iceberg}{non-traded goods} $(z_E, z_M)$, bought only at Home. Its wage bill equals
spending on its goods:
$wL = \vartheta(z_E)(wL + w^*L^*) + (\vartheta(z_M) - \vartheta(z_E))wL$, where \omterm{def:b3:ricardian-trade:iceberg}{iceberg}
losses are paid to Home's workers as part of the delivered price. Simplifying,
$wL(1 - \vartheta(z_M)) = \vartheta(z_E)w^*L^*$, which is the stated equation. As $\omega$
rises, $z_E$ and $z_M$ fall because $A$ is decreasing, so $\vartheta(z_E)$ falls and $1 -
\vartheta(z_M)$ rises: the left side increases strictly with $\omega$ and the right side
decreases. The left side is near 0 for small $\omega$ (where $z_M$ is near 1) and the
right side near 0 for large $\omega$, so they cross exactly once.
\end{solution}

\begin{solution}{exo:b3:ricardian-trade:11}
At the old $\omega$, $\kappa\omega$ falls and $\omega/\kappa$ rises: goods at the top of the band
($A$ just below $1.2\omega$) become Home exports and goods at its bottom ($A$ just
above $\omega/1.2$) Home imports. With equal spending $\Delta$ on each end, the right
side of the equation of \cref{exo:b3:ricardian-trade:10} rises by $\Delta L^*$ and the
left by $\omega\Delta L$. With $L = L^*$, the left rises by more if $\omega > 1$: Home's new
imports cost more than its new exports earn, and $\omega$ must fall; if $\omega < 1$, it
rises.
\end{solution}

\begin{solution}{exo:b3:ricardian-trade:12}
Inside the model: the marginal good, at $A = \omega$, is made in both countries;
trade costs leave a band of goods made in both; an industry in the statistics
groups many goods with different $A_i$, so an industry can both export and
import. Outside it: products are differentiated by firm, and countries trade
varieties of the same good (\cref{ch:b3:scale-variety-firms}); other factors
(land, capital, skills) and rising costs make specialisation incomplete
(\cref{ch:b3:factor-endowments}).
\end{solution}

\begin{solution}{pb:b3:ricardian-trade:twelve}
\textbf{1.} $A$: grain $12/2 = 6$, timber 4, fish 3, wool 2.4, cheese 2, paper 1.75,
cloth 1.5, glass 1.38, tools 1.25, shoes 1, clocks 0.8, radios 0.6: decreasing down
the table. % figdata: twelve_goods

\textbf{2.} Brenland needs fewer hours in the first nine goods, as many in shoes,
more in clocks and radios. It does not matter: the pattern depends on relative
costs, and the wage absorbs absolute differences.

\textbf{3.} Cloth in grain: Brenland $4/2 = 2$, Kestria $6/12 = 0.5$. Radios in grain:
Brenland $10/2 = 5$, Kestria $6/12 = 0.5$.

\textbf{4.} At $\omega = 1$: the nine goods with $A > 1$, with shoes ($A = 1$) made in
both. At $\omega = 2$: grain, timber, fish and wool, with cheese made in both.

\textbf{5.} Above 6 Brenland would make nothing and below 0.6 everything; neither
lets trade balance.

\textbf{6.} If Brenland makes $e$ goods and Kestria $m = 12 - e$, world spending on
Brenland's goods is $e/12$ of world income: $wL = \tfrac{e}{12}(wL + w^*L)$, so $\omega(12 -
e) = e$ and $\omega = e/m$.

\textbf{7.} $7/5 = 1.4$, and $A_7 = 1.5 > 1.4 > 1.38 = A_8$: at that wage Brenland makes
exactly the first seven goods. % figdata: twelve_results

\textbf{8.} At 1.2 Brenland would make the nine goods with $A > 1.2$, and trade would
balance only at $9/3 = 3$. At 2 it would make four goods and part of cheese,
which balances at a wage between $4/8 = 0.5$ and $5/7 = 0.71$, not 2.

\textbf{9.} Autarky: a unit of grain buys $2/10 = 0.2$ radios. With trade: $p_{\text{grain}}
= 2w$, $p_{\text{radio}} = 6w^*$, so grain buys $2\omega/6 = 2.8/6 = 0.47$ radios.

\textbf{10.} Autarky: $1/10 = 0.1$ radio an hour. Trade: $w/(6w^*) = 1.4/6 = 0.23$.

\textbf{11.} With equal shares the price index is the geometric mean of the twelve
prices. Goods Brenland makes keep their price $wa_i$; an imported good's price
falls from $wa_i$ to $w^*a_i^*$, a ratio of $a_i\omega/a_i^* = \omega/A_i$, so real income is
multiplied by the twelfth root of the product of those ratios: $1.0145 \times 1.12
\times 1.4 \times 1.75 \times 2.333 = 6.496$, whose twelfth root is 1.169: a gain of
16.9~\%.

\textbf{12.} Kestria imports the first seven goods; the ratios $A_i/\omega$ are 4.286,
2.857, 2.143, 1.714, 1.429, 1.25 and 1.071, with product 86.06 and twelfth root 1.450:
45.0~\%. Kestria gains more because world prices, set by the relative wage of 1.4,
are far from its \omterm{def:b3:ricardian-trade:autarky}{autarky prices} for the goods it imports (grain cost it six
times as much labour relative to Brenland), while for Brenland's imports the
ratios are smaller.

\textbf{13.} Export if $A_i > 1.1\omega$; import if $A_i < \omega/1.1$; make for itself
between the two.

\textbf{14.} $1.1 \times 1.5 = 1.65$ and $1.5/1.1 = 1.364$. Exports: grain, timber, fish,
wool, cheese, paper (6). Non-traded: cloth, glass. Imports: tools, shoes, clocks,
radios (4).

\textbf{15.} Brenland's exports earn Kestria's spending on six goods, $\tfrac{6}{12}w^*L$;
its imports cost $\tfrac{4}{12}wL$. Balance requires $4\omega = 6$, $\omega = 1.5$, and the
groups of Question 14 are those of that wage.

\textbf{16.} Cloth leaves Brenland's exports and glass its imports. Losing one
import out of five cuts Brenland's spending abroad by a larger proportion than
losing one export out of seven cuts its earnings, so its wage must rise to
restore balance: $6/4 > 7/5$.

\textbf{17.} The ratios $\omega/(\kappa A_i)$ for the four imports are 1.091, 1.364, 1.705
and 2.273, with product 5.763 and twelfth root 1.157: 15.7~\%. Both countries gain
less than in free trade (16.9~\% and 45.0~\%); Kestria loses more.

\textbf{18.} $A_i/1.2$: 5, 3.333, 2.5, 2, 1.667, 1.458, 1.25, 1.15, 1.042, 0.833, 0.667, 0.5.

\textbf{19.} If Brenland made the six goods with $A > 1.25$, trade would balance at
$6/6 = 1$, below $A_{\text{cloth}} = 1.25$; if it made seven, at $7/5 = 1.4$, above it.
So cloth is shared and $\omega = 1.25$. Brenland's share of world income is then
$1.25/2.25 = 0.556$, which must equal $(6 + s)/12$: $s = 0.667$, two thirds of the
world's cloth.

\textbf{20.} 20.3~\% against 16.9~\%: Brenland's real wage rises, as
\cref{prop:b3:ricardian-trade:abroad} says. % figdata: twelve_results

\textbf{21.} No. Its autarky real wage has risen by 20~\% in every good, and its
gain over that new autarky is 39.2~\%: its real wage is $1.2 \times 1.392 = 1.67$ times
its old autarky level, against 1.45 before.

\textbf{22.} Brenland spends on the seven goods it makes: $\lambda = 7/12 = 0.583$. The
formula gives $0.583^{-1/5} - 1 = 11.4~\%$.

\textbf{23.} $\varepsilon = \ln(12/7)/\ln 1.169 = 0.539/0.156 = 3.46$.

\textbf{24.} The formula holds when the trade elasticity is constant, as with
the productivity distributions of Eaton and Kortum. With twelve goods and fixed
labour requirements, trade responds to costs in jumps, as whole goods switch
country, and no single elasticity describes it.

\textbf{25.} Computed directly from the prices of imported goods, Brenland's gain
from free trade is 16.9~\% of real income; the \omterm{def:b3:ricardian-trade:acr}{gains-from-trade formula} with an
elasticity of 5 gives 11.4~\%, and the two agree only with an elasticity of 3.46.
\end{solution}
